\documentclass[../main.tex]{subfiles}

\begin{document}

\chapter{fake-reco} \label{ch:fake-reco}
  fake-reco è un processo il cui scopo è copiare la verità Monte Carlo
  scritta sui file generati da GENIE ed EDep-Sim nei file di output di
  SANDReco (i CAF), sia nei campi dedicati alla verità stessa, sia in quelli
  dedicati alla ricostruzione; per fare ciò deve tradurre in maniera
  opportuna le strutture dati presenti negli output di EDep-Sim, riportate
  in figura~\ref{fig:input-structure}, in quelle definite dai CAF\@.\\

  \subfile{obbiettivi/presenti.tex}

  \subfile{obbiettivi/futuri.tex}

  \subfile{implementazione/implementazione.tex}


%    Come illustrato i figura~\ref{fig:fake-reco-v0} fake-reco è implementato
%    nel ufw come processo e attualmente non accetta alcun parametro di
%    configurazione, quando sarà aggiunta la possibilità di applicare uno
%    \textit{smearing} sarà introdotto un parametro contestuale.\\


    \subfile{json.tex}

\end{document}
